% Part: sets-functions-relations
% Chapter: relations
% Section: orders
% Hindi translation (hi-Deva-IN) of Open Logic commit 9620cc73f9c8e0ad003c514a5d3748f29611c4c0.

\documentclass[../../../include/open-logic-section]{subfiles}

\begin{document}

\olfileid[hi]{sfr}{rel}{ord}
\olsection{क्रम}

\begin{explain}
हमारी अनेक तुलनाओं में किसी निश्चित दृष्टि से कुछ वस्तुओं को अन्य वस्तुओं से
“छोटा”, “बराबर”, या “बड़ा” बताया जाता है। इनमें \emph{क्रम} संबंध आते
हैं। परंतु क्रम संबंध कई प्रकार के होते हैं। उदाहरण के लिए, कुछ में किन्हीं
भी दो वस्तुओं का तुलनीय होना आवश्यक होता है, जबकि अन्य में नहीं। कुछ बराबरी
को सम्मिलित करते हैं (जैसे~$\le$) और कुछ उसे अलग रखते हैं (जैसे~$<$)। यहाँ
इनका वर्गीकरण उपयोगी होगा।
\end{explain}

\begin{defn}[पूर्वक्रम]
जो संबंध स्वतुल्य और संक्रामक दोनों हो, उसे
\emph{पूर्वक्रम} कहते हैं।  
\end{defn}

\begin{defn}[आंशिक क्रम]
जो पूर्वक्रम प्रति-सममित भी हो, उसे
\emph{आंशिक क्रम} कहते हैं।
\end{defn}

\begin{defn}[रैखिक क्रम]\ollabel{def:linearorder}
जो आंशिक क्रम संबद्ध भी हो, उसे
\emph{संपूर्ण क्रम} या \emph{रैखिक क्रम} कहते हैं।
\end{defn}

\begin{ex}
हर रैखिक क्रम एक आंशिक क्रम भी होता है, और हर आंशिक क्रम एक पूर्वक्रम भी
होता है, परंतु इनके उलटे कथन सत्य नहीं हैं।~$A$ पर सार्वत्रिक संबंध एक
पूर्वक्रम है, क्योंकि वह स्वतुल्य और संक्रामक है। किंतु यदि
$A$ में एक से अधिक !!{element} हों, तो सार्वत्रिक संबंध
प्रति-सममित नहीं होता और इसलिए आंशिक क्रम भी नहीं होता।
\end{ex}

\begin{ex}
\emph{इससे अधिक लंबा नहीं है} संबंध $\preccurlyeq$ पर, जो~$\Bin^*$ पर है,
विचार कीजिए: $x \preccurlyeq y$ तभी और केवल तभी जब $\len{x} \le \len{y}$। यह एक
पूर्वक्रम (स्वतुल्य और संक्रामक) है, और संबद्ध भी है, परंतु आंशिक क्रम नहीं
है, क्योंकि यह प्रति-सममित नहीं है। उदाहरण के लिए, $01
\preccurlyeq 10$ और $10 \preccurlyeq 01$, परंतु $01 \neq 10$।
\end{ex}

\begin{ex}
समुच्चयों के किसी समुच्चय पर संबंध $\subseteq$ एक महत्त्वपूर्ण आंशिक क्रम
है। सामान्यतः यह रैखिक क्रम नहीं है, क्योंकि यदि $a \neq b$ हो और हम
$\Pow{\{a, b\}} = \{\emptyset, \{a\}, \{b\}, \{a,b\}\}$ पर विचार करें, तो
दिखता है कि $\{a\} \nsubseteq \{b\}$ और $\{a\} \neq \{b\}$ और $\{b\}
\nsubseteq \{a\}$।
\end{ex}

\begin{ex}
\emph{निःशेष विभाज्यता} का संबंध हमें ऐसा आंशिक क्रम देता है जो रैखिक क्रम
नहीं है। पूर्णांकों $n$ और~$m$ के लिए हम $n \mid m$ लिखकर यह बताते हैं कि
$n$, $m$ को निःशेष विभाजित करता है, अर्थात तभी और केवल तभी जब कोई
पूर्णांक~$k$ ऐसा हो कि $m = kn$।~$\Nat$ पर यह एक आंशिक क्रम है, परंतु रैखिक
क्रम नहीं है: उदाहरण के लिए, $2 \nmid 3$ और $3 \nmid
2$ भी। $\Int$ पर संबंध के रूप में देखने पर विभाज्यता केवल पूर्वक्रम है,
क्योंकि यह प्रति-सममित नहीं है: $1 \mid -1$ और $-1 \mid 1$
परंतु $1 \neq -1$।
\end{ex}

\begin{ex}
अनुक्रमों के समुच्चय~$A^*$ पर \emph{विस्तार} संबंध यह है: $s \sqsubseteq s'$
तभी और केवल तभी जब $s = \emptyseq$ (रिक्त
अनुक्रम), $s = s'$, अथवा $s = \tuple{s_1, \dots, s_n}$ और $s' =
\tuple{s_1, \dots, s_n, s_{n+1}, \dots, s_m}$। यदि $s \sqsubseteq s'$ हो,
तो हम यह भी कहते हैं कि $s$,~$s'$ का \emph{आरंभिक खंड} है।~$A^*$ पर
विस्तार संबंध एक आंशिक क्रम है, परंतु रैखिक क्रम नहीं है; उदाहरणार्थ, यदि
$a \neq b$, तो $ab \not\sqsubseteq ba$ और $ba \not\sqsubseteq ab$।
\end{ex}

\begin{defn}[प्रबल क्रम]
\emph{प्रबल क्रम} ऐसा संबंध है जो अस्वतुल्य, असममित और संक्रामक हो।
\end{defn}

\begin{defn}[प्रबल रैखिक क्रम]\ollabel{def:strictlinearorder}
जो प्रबल क्रम संबद्ध भी हो, उसे 
\emph{प्रबल संपूर्ण क्रम} या \emph{प्रबल रैखिक क्रम} कहते हैं।
\end{defn}

\begin{ex}
$\le$ वह रैखिक क्रम है जो प्रबल रैखिक क्रम~$<$ के अनुरूप है। $\subseteq$
वह आंशिक क्रम है जो प्रबल क्रम~$\subsetneq$ के अनुरूप है।
\end{ex}

किसी भी प्रबल क्रम $R$ में, जो~$A$ पर है, विकर्ण $\Id{A}$ जोड़कर, अर्थात
सभी युग्म~$\tuple{x,
x}$ जोड़कर, उसे आंशिक क्रम में बदला जा सकता है। (इसे
~$R$ का \emph{स्वतुल्य संवरक} कहते हैं।) इसके विपरीत, किसी आंशिक क्रम से
आरंभ करके~$\Id{A}$ हटाने पर प्रबल क्रम प्राप्त किया जा सकता है। अगले दो
परिणाम इसे सटीक रूप देते हैं।

\begin{prop}\ollabel{prop:stricttopartial}
यदि $R$,~$A$ पर एक प्रबल क्रम है, तो $R^+ = R \cup \Id{A}$ एक आंशिक क्रम
है। इसके अतिरिक्त, यदि $R$ एक प्रबल रैखिक क्रम है, तो $R^+$ एक रैखिक क्रम
है।
\end{prop}

\begin{proof}
मान लें कि $R$ एक प्रबल क्रम है, अर्थात $R \subseteq A^2$ और $R$
अस्वतुल्य, असममित और संक्रामक है। मान लें $R^+ = R \cup \Id{A}$। हमें
दिखाना है कि $R^+$ स्वतुल्य, प्रति-सममित और संक्रामक है।

$R^+$ स्पष्टतः स्वतुल्य है, क्योंकि $\tuple{x, x} \in \Id{A} \subseteq
R^+$ प्रत्येक $x \in A$ के लिए। 

$R^+$ को प्रति-सममित दिखाने के लिए, विरोधाभासार्थ मान लें कि $R^+xy$ और
$R^+yx$, परंतु $x \neq y$। चूँकि $\tuple{x,y} \in R \cup \Id{A}$, परंतु
$\tuple{x, y} \notin \Id{A}$, इसलिए $\tuple{x, y} \in R$ होना ही चाहिए,
अर्थात~$Rxy$। इसी प्रकार,~$Ryx$। किंतु यह हमारी उस धारणा का खंडन करता है
कि $R$ असममित है।

संक्रामकता स्थापित करने के लिए मान लें कि $R^+xy$ और $R^+yz$। यदि
$\tuple{x, y} \in R$ और $\tuple{y,z} \in R$ दोनों हों, तो $\tuple{x, z} \in
R$, क्योंकि $R$~संक्रामक है। अन्यथा, या तो $\tuple{x, y} \in
\Id{A}$, अर्थात $x = y$, अथवा $\tuple{y, z} \in \Id{A}$, अर्थात $y = z$।
पहली स्थिति में धारणा से $R^+yz$ है और $x = y$, इसलिए
$R^+xz$। दूसरी स्थिति में भी ठीक इसी प्रकार। दोनों स्थितियों में $R^+xz$;
अतः $R^+$ भी संक्रामक है।

“इसके अतिरिक्त” वाले अंश के लिए मान लें कि $R$ संबद्ध भी है। अतः सभी
$x \neq y$ के लिए या तो $Rxy$ अथवा~$Ryx$, अर्थात या तो
$\tuple{x, y} \in R$ अथवा $\tuple{y, x} \in R$। चूँकि $R \subseteq R^+$,
यह बात $R^+$ के लिए भी सत्य रहती है, इसलिए $R^+$ भी संबद्ध है।
\end{proof}

\begin{prop}\ollabel{prop:partialtostrict}
यदि $R$,~$A$ पर एक आंशिक क्रम है, तो $R^- = R \setminus \Id{A}$ एक प्रबल
क्रम है। इसके अतिरिक्त, यदि $R$ एक रैखिक क्रम है, तो $R^-$ एक प्रबल
रैखिक क्रम है।
\end{prop}

\begin{proof}
इसे अभ्यास के लिए छोड़ा जाता है।
\end{proof}

\begin{prob}
\olref[sfr][rel][ord]{prop:partialtostrict} का प्रमाण दीजिए। 
\end{prob}

निम्न सरल परिणाम स्थापित करता है कि प्रबल रैखिक क्रमों में
विस्तारिता-सदृश गुणधर्म होता है:

\begin{prop}\ollabel{prop:extensionality-strictlinearorders}
यदि $<$, $A$ पर एक प्रबल रैखिक क्रम है, तो:
\[
  (\forall a, b \in A)((\forall x \in A)(x < a \liff x < b) \lif a = b).
\]
\end{prop}

\begin{proof}
मान लें $(\forall x \in A)(x < a \liff x < b)$। यदि $a < b$, तो $a <
a$, जो $<$ के अस्वतुल्य होने का खंडन करता है; इसलिए $a \nless b$।
ठीक इसी प्रकार, $b \nless a$। अतः $a = b$, क्योंकि $<$ संबद्ध है।
\end{proof}

\end{document}
